\documentclass[11pt, a4paper]{article}
\usepackage[a4paper, top=2.5cm, bottom=2.5cm, left=2cm, right=2cm]{geometry}
\usepackage{fontspec}
\usepackage[english, bidi=basic, provide=*]{babel}
\babelprovide[import, onchar=ids fonts]{english}
\babelfont{rm}{TeX Gyre Termes}
\babelfont{sf}{TeX Gyre Heros}
\babelfont{tt}{TeX Gyre Cursor}
\usepackage{enumitem}
\usepackage{amsmath}

\title{MCQ Practice Set: Mendel and the Gene Idea}
\author{Subject: Biology (Chapter 14)}
\date{}

\begin{document}
\maketitle

\section*{Practice Questions}

\begin{enumerate}[leftmargin=*, label=\arabic*.]

\item \textbf{Mendel's "hereditary factors" are now known as:}
\begin{enumerate}[label=\Alph*.]
\item Chromosomes
\item Proteins
\item Genes
\item Alleles
\item Gametes
\end{enumerate}
\textbf{Answer:} C \\
\textbf{Explanation:} What Mendel called "factors" acting as discrete units of inheritance are today identified as genes.

\item \textbf{Which of the following is the reason Mendel chose pea plants for his experiments?}
\begin{enumerate}[label=\Alph*.]
\item They have short generation times.
\item They produce a large number of offspring.
\item Mating could be strictly controlled.
\item They have many varieties with distinct heritable features.
\item All of the above.
\end{enumerate}
\textbf{Answer:} E \\
\textbf{Explanation:} Peas provided the perfect model system due to their rapid reproduction, high yield, controllable pollination, and clear-cut traits.

\item \textbf{A heritable feature that varies among individuals, such as flower color, is called a \underline{\hspace{1cm}}, while each variant is called a \underline{\hspace{1cm}}.}
\begin{enumerate}[label=\Alph*.]
\item Trait; character
\item Character; trait
\item Allele; gene
\item Genotype; phenotype
\item Locus; allele
\end{enumerate}
\textbf{Answer:} B \\
\textbf{Explanation:} A character is a functional feature (like color), and a trait is the specific version of that character (like purple or white).

\item \textbf{In Mendel's experiments, the "true-breeding" parents are referred to as the \underline{\hspace{1cm}} generation.}
\begin{enumerate}[label=\Alph*.]
\item F1
\item F2
\item P
\item G
\item Hybrid
\end{enumerate}
\textbf{Answer:} C \\
\textbf{Explanation:} P stands for Parental; true-breeding means that self-pollination always produces offspring with the same trait.

\item \textbf{When Mendel crossed true-breeding purple-flowered plants with true-breeding white-flowered plants, the F1 generation:}
\begin{enumerate}[label=\Alph*.]
\item Was all white
\item Was all light purple
\item Was all purple
\item Was a 3:1 mix of purple and white
\item Was sterile
\end{enumerate}
\textbf{Answer:} C \\
\textbf{Explanation:} The purple trait was dominant, completely masking the white trait in the F1 hybrids.

\item \textbf{What phenotypic ratio did Mendel consistently observe in the F2 generation of a monohybrid cross?}
\begin{enumerate}[label=\Alph*.]
\item 1:1
\item 3:1
\item 1:2:1
\item 9:3:3:1
\item 2:1
\end{enumerate}
\textbf{Answer:} B \\
\textbf{Explanation:} The dominant trait appeared in three-quarters of the offspring, while the recessive trait reappeared in one-quarter.

\item \textbf{Mendel’s Model explains that alternative versions of genes account for variations in inherited characters. These versions are called:}
\begin{enumerate}[label=\Alph*.]
\item Loci
\item Chromosomes
\item Alleles
\item Traits
\item Genotypes
\end{enumerate}
\textbf{Answer:} C \\
\textbf{Explanation:} Alleles are different molecular sequences of the same gene located at the same locus on homologous chromosomes.

\item \textbf{The law of segregation states that:}
\begin{enumerate}[label=\Alph*.]
\item Genes for different traits sort independently.
\item Alleles of a gene are always identical.
\item The two alleles for a heritable character separate during gamete formation.
\item Dominant alleles are more common in a population than recessive ones.
\item Offspring always resemble the mother.
\end{enumerate}
\textbf{Answer:} C \\
\textbf{Explanation:} Segregation occurs during meiosis, ensuring each gamete receives only one of the two alleles present in the somatic cells.

\item \textbf{An organism with two identical alleles for a character is \underline{\hspace{1cm}}, while an organism with two different alleles is \underline{\hspace{1cm}}.}
\begin{enumerate}[label=\Alph*.]
\item Heterozygous; homozygous
\item Phenotypic; genotypic
\item Homozygous; heterozygous
\item Dominant; recessive
\item Recessive; dominant
\end{enumerate}
\textbf{Answer:} C \\
\textbf{Explanation:} "Homo" means same; "Hetero" means different.

\item \textbf{Which term refers to an organism's physical appearance, as opposed to its genetic makeup?}
\begin{enumerate}[label=\Alph*.]
\item Genotype
\item Phenotype
\item Allele frequency
\item Haplotype
\item Pedigree
\end{enumerate}
\textbf{Answer:} B \\
\textbf{Explanation:} Phenotype is the observable trait; genotype is the underlying combination of alleles.

\item \textbf{To determine the genotype of an individual with a dominant phenotype, a geneticist would cross it with a homozygous recessive individual. This is called a:}
\begin{enumerate}[label=\Alph*.]
\item Dihybrid cross
\item Monohybrid cross
\item Testcross
\item Backcross
\item Punnett square
\end{enumerate}
\textbf{Answer:} C \\
\textbf{Explanation:} A testcross reveals if the dominant individual is homozygous or heterozygous based on the offspring's phenotypes.

\item \textbf{Mendel’s second law, the Law of Independent Assortment, was derived by following \underline{\hspace{1cm}} characters at the same time.}
\begin{enumerate}[label=\Alph*.]
\item One
\item Two
\item Three
\item All seven
\item Linked
\end{enumerate}
\textbf{Answer:} B \\
\textbf{Explanation:} Dihybrid crosses showed that the inheritance of one character has no effect on the inheritance of another (if the genes are on different chromosomes).

\item \textbf{What is the expected phenotypic ratio in the F2 generation of a dihybrid cross ($YyRr \times YyRr$)?}
\begin{enumerate}[label=\Alph*.]
\item 3:1
\item 1:2:1
\item 9:3:3:1
\item 1:1:1:1
\item 9:7
\end{enumerate}
\textbf{Answer:} C \\
\textbf{Explanation:} The 9:3:3:1 ratio occurs when two independent traits (each a 3:1 ratio) are combined ($3 \times 3 : 3 \times 1 : 1 \times 3 : 1 \times 1$).

\item \textbf{The multiplication rule of probability is used to determine the probability that:}
\begin{enumerate}[label=\Alph*.]
\item Either of two independent events will occur.
\item Two or more independent events will occur together in some specific combination.
\item An event will not occur.
\item A dominant trait will appear.
\item A mutation will happen.
\end{enumerate}
\textbf{Answer:} B \\
\textbf{Explanation:} To find the probability of events A AND B, you multiply their individual probabilities.

\item \textbf{If you toss a coin twice, what is the probability that it will land on heads both times?}
\begin{enumerate}[label=\Alph*.]
\item 1/2
\item 1/4
\item 1/8
\item 3/4
\item 1
\end{enumerate}
\textbf{Answer:} B \\
\textbf{Explanation:} $1/2 \times 1/2 = 1/4$.

\item \textbf{The addition rule of probability states that the probability of any one of two or more exclusive events occurring is calculated by:}
\begin{enumerate}[label=\Alph*.]
\item Multiplying their individual probabilities.
\item Subtracting the smaller from the larger probability.
\item Adding their individual probabilities together.
\item Dividing one by the other.
\item Taking the square root.
\end{enumerate}
\textbf{Answer:} C \\
\textbf{Explanation:} The addition rule is used for "OR" scenarios (e.g., getting a heterozygous genotype as $Pp$ or $pP$).

\item \textbf{In a cross $AaBbCc \times AaBbCc$, what is the probability of producing an offspring with the genotype $aabbcc$?}
\begin{enumerate}[label=\Alph*.]
\item 1/4
\item 1/16
\item 1/64
\item 3/64
\item 1/8
\end{enumerate}
\textbf{Answer:} C \\
\textbf{Explanation:} Prob($aa$) = 1/4; Prob($bb$) = 1/4; Prob($cc$) = 1/4. Multiply them: $1/4 \times 1/4 \times 1/4 = 1/64$.

\item \textbf{When the phenotype of the F1 hybrid is somewhere between the phenotypes of the two parental varieties, it is called:}
\begin{enumerate}[label=\Alph*.]
\item Complete dominance
\item Codominance
\item Incomplete dominance
\item Pleiotropy
\item Epistasis
\end{enumerate}
\textbf{Answer:} C \\
\textbf{Explanation:} An example is snapdragons, where red and white parents produce pink F1 offspring.

\item \textbf{Human blood groups (ABO) are an example of \underline{\hspace{1cm}} and \underline{\hspace{1cm}}.}
\begin{enumerate}[label=\Alph*.]
\item Complete dominance; epistasis
\item Incomplete dominance; pleiotropy
\item Multiple alleles; codominance
\item Polygenic inheritance; incomplete dominance
\item Lethal alleles; sex-linkage
\end{enumerate}
\textbf{Answer:} C \\
\textbf{Explanation:} There are three alleles ($I^A, I^B, i$), and $I^A$ and $I^B$ are codominant.

\item \textbf{A single gene that has multiple phenotypic effects, such as the gene causing cystic fibrosis, is an example of:}
\begin{enumerate}[label=\Alph*.]
\item Epistasis
\item Polygenic inheritance
\item Pleiotropy
\item Blending
\item Incomplete dominance
\end{enumerate}
\textbf{Answer:} C \\
\textbf{Explanation:} Pleiotropy occurs when one gene influences several seemingly unrelated traits.

\item \textbf{In \underline{\hspace{1cm}}, the phenotypic expression of a gene at one locus alters that of a gene at a second locus.}
\begin{enumerate}[label=\Alph*.]
\item Codominance
\item Pleiotropy
\item Epistasis
\item Incomplete dominance
\item Independent assortment
\end{enumerate}
\textbf{Answer:} C \\
\textbf{Explanation:} A classic example is coat color in Labradors, where a second gene determines if pigment is deposited at all.

\item \textbf{Skin color in humans, which varies along a continuum, is an example of:}
\begin{enumerate}[label=\Alph*.]
\item Pleiotropy
\item Epistasis
\item Polygenic inheritance
\item Codominance
\item Complete dominance
\end{enumerate}
\textbf{Answer:} C \\
\textbf{Explanation:} Polygenic inheritance involves the additive effect of two or more genes on a single phenotypic character.

\item \textbf{Characters that are influenced by both genetic and environmental factors are called:}
\begin{enumerate}[label=\Alph*.]
\item Quantitative
\item Multifactorial
\item Mendelian
\item Heterozygous
\item Lethal
\end{enumerate}
\textbf{Answer:} B \\
\textbf{Explanation:} Multifactorial characters emphasize the collective influence of nature and nurture.

\item \textbf{A diagram showing the occurrence and appearance of phenotypes of a particular gene or organism and its ancestors from one generation to the next is a:}
\begin{enumerate}[label=\Alph*.]
\item Punnett square
\item Karyotype
\item Pedigree
\item Genome map
\item Flowchart
\end{enumerate}
\textbf{Answer:} C \\
\textbf{Explanation:} Pedigrees are used to trace inheritance patterns in families, especially for human genetic disorders.

\item \textbf{Heterozygotes who are phenotypically normal but can transmit a recessive allele to their offspring are called:}
\begin{enumerate}[label=\Alph*.]
\item Hybrids
\item Mutants
\item Carriers
\item Vectors
\item Progenitors
\end{enumerate}
\textbf{Answer:} C \\
\textbf{Explanation:} Carriers have one recessive allele for a disorder but do not show symptoms because the dominant allele is functional.

\item \textbf{Which of the following is a common autosomal recessive disorder?}
\begin{enumerate}[label=\Alph*.]
\item Huntington's disease
\item Achondroplasia
\item Cystic fibrosis
\item Polydactyly
\item Hypercholesterolemia
\end{enumerate}
\textbf{Answer:} C \\
\textbf{Explanation:} Cystic fibrosis is the most common lethal genetic disease in the United States among people of European descent.

\item \textbf{Sickle-cell disease is caused by the substitution of a single amino acid in the \underline{\hspace{1cm}} protein of red blood cells.}
\begin{enumerate}[label=\Alph*.]
\item Insulin
\item Hemoglobin
\item Keratin
\item Collagen
\item Actin
\end{enumerate}
\textbf{Answer:} B \\
\textbf{Explanation:} This substitution causes hemoglobin to aggregate, deforming the red blood cells into a sickle shape.

\item \textbf{What is the evolutionary advantage of being a carrier for the sickle-cell allele in some regions?}
\begin{enumerate}[label=\Alph*.]
\item Resistance to tuberculosis
\item Faster blood clotting
\item Resistance to malaria
\item Enhanced oxygen transport
\item Immunity to the flu
\end{enumerate}
\textbf{Answer:} C \\
\textbf{Explanation:} Heterozygotes have increased resistance to the malaria parasite, a phenomenon called heterozygote advantage.

\item \textbf{Huntington's disease is unique because it is caused by a \underline{\hspace{1cm}} allele that does not manifest symptoms until the individual is \underline{\hspace{1cm}}.}
\begin{enumerate}[label=\Alph*.]
\item Recessive; newborn
\item Dominant; 35-45 years old
\item Codominant; a teenager
\item Lethal; in utero
\item Pleiotropic; elderly
\end{enumerate}
\textbf{Answer:} B \\
\textbf{Explanation:} Because it is dominant and late-onset, the allele is often passed to offspring before the parent knows they are affected.

\item \textbf{What is the probability of a child being born with a recessive condition if both parents are carriers?}
\begin{enumerate}[label=\Alph*.]
\item 0\%
\item 25\%
\item 50\%
\item 75\%
\item 100\%
\end{enumerate}
\textbf{Answer:} B \\
\textbf{Explanation:} $Aa \times Aa$ produces 1/4 $AA$, 1/2 $Aa$, and 1/4 $aa$.

\item \textbf{Which technique involves sampling the liquid that bathes the fetus to test for genetic disorders?}
\begin{enumerate}[label=\Alph*.]
\item Chorionic villus sampling (CVS)
\item Amniocentesis
\item Ultrasound
\item Fetoscopy
\item IVF
\end{enumerate}
\textbf{Answer:} B \\
\textbf{Explanation:} Amniocentesis collects amniotic fluid containing fetal cells for biochemical and genetic testing.

\item \textbf{In a cross $Pp \times Pp$, what proportion of the purple-flowered offspring are expected to be true-breeding?}
\begin{enumerate}[label=\Alph*.]
\item 1/4
\item 1/2
\item 1/3
\item 2/3
\item 3/4
\end{enumerate}
\textbf{Answer:} C \\
\textbf{Explanation:} Of the 3/4 purple offspring ($AA, Aa, aA$), only 1/4 ($AA$) are true-breeding. $1/4 \div 3/4 = 1/3$.

\item \textbf{Mendel’s "Law of Independent Assortment" applies only to genes that are:}
\begin{enumerate}[label=\Alph*.]
\item On the same chromosome
\item Located near each other
\item On different chromosomes or far apart on the same chromosome
\item Responsible for the same trait
\item Always recessive
\end{enumerate}
\textbf{Answer:} C \\
\textbf{Explanation:} Linked genes on the same chromosome tend to be inherited together and do not assort independently.

\item \textbf{If a plant with genotype $YyRr$ produces gametes, which of the following is NOT a possible gamete (assuming independent assortment)?}
\begin{enumerate}[label=\Alph*.]
\item $YR$
\item $Yr$
\item $yR$
\item $yr$
\item $Yy$
\end{enumerate}
\textbf{Answer:} E \\
\textbf{Explanation:} Gametes are haploid and must contain one allele of each gene ($Y$ or $y$ AND $R$ or $r$).

\item \textbf{An allele that is only expressed in the phenotype when the dominant allele is absent is:}
\begin{enumerate}[label=\Alph*.]
\item Homozygous
\item Recessive
\item Codominant
\item Pleiotropic
\item Epistatic
\end{enumerate}
\textbf{Answer:} B \\
\textbf{Explanation:} Recessive traits are masked by dominant alleles in heterozygotes.

\item \textbf{The multiplication of probabilities for independent events is based on the word \underline{\hspace{1cm}}, while addition is based on the word \underline{\hspace{1cm}}.}
\begin{enumerate}[label=\Alph*.]
\item Or; and
\item And; or
\item If; then
\item Always; never
\item Both; neither
\end{enumerate}
\textbf{Answer:} B \\
\textbf{Explanation:} Probability(A AND B) = mult; Probability(A OR B) = add.

\item \textbf{Mendel’s F1 hybrids were always \underline{\hspace{1cm}} for the trait being studied.}
\begin{enumerate}[label=\Alph*.]
\item Homozygous dominant
\item Homozygous recessive
\item Heterozygous
\item Mutants
\item Haploid
\end{enumerate}
\textbf{Answer:} C \\
\textbf{Explanation:} Crossing two different true-breeding parents ($AA \times aa$) produces heterozygous offspring ($Aa$).

\item \textbf{The phenotype of a $I^A I^B$ individual is blood type AB. This is an example of:}
\begin{enumerate}[label=\Alph*.]
\item Incomplete dominance
\item Codominance
\item Epistasis
\item Pleiotropy
\item Polygenic inheritance
\end{enumerate}
\textbf{Answer:} B \\
\textbf{Explanation:} Both alleles are fully expressed in the phenotype.

\item \textbf{A testcross of an individual with an unknown genotype results in a 1:1 ratio of dominant to recessive phenotypes. The unknown parent was:}
\begin{enumerate}[label=\Alph*.]
\item Homozygous dominant
\item Heterozygous
\item Homozygous recessive
\item A mutation
\item Sterile
\end{enumerate}
\textbf{Answer:} B \\
\textbf{Explanation:} $Aa \times aa$ yields 1/2 $Aa$ and 1/2 $aa$.

\item \textbf{How many different gamete types can be produced by an individual with the genotype $AaBbCcDd$?}
\begin{enumerate}[label=\Alph*.]
\item 4
\item 8
\item 16
\item 32
\item 64
\end{enumerate}
\textbf{Answer:} C \\
\textbf{Explanation:} The formula is $2^n$, where $n$ is the number of heterozygous pairs. $2^4 = 16$.

\item \textbf{Which of the following is true about dominant alleles?}
\begin{enumerate}[label=\Alph*.]
\item they are always more common in the population.
\item They are "stronger" and physically subdue recessive alleles.
\item They are expressed in the phenotype even if only one copy is present.
\item They are always more advantageous for survival.
\item They prevent the recessive allele from being transcribed.
\end{enumerate}
\textbf{Answer:} C \\
\textbf{Explanation:} Dominance is a description of the relationship between alleles in the phenotype, not a measure of frequency or strength.

\item \textbf{Tay-Sachs disease is \underline{\hspace{1cm}} at the organismal level, \underline{\hspace{1cm}} at the biochemical level, and \underline{\hspace{1cm}} at the molecular level.}
\begin{enumerate}[label=\Alph*.]
\item Dominant; Recessive; Codominant
\item Recessive; Incompletely dominant; Codominant
\item Recessive; Dominant; Incompletely dominant
\item Codominant; Recessive; Incompletely dominant
\item Recessive; Codominant; Dominant
\end{enumerate}
\textbf{Answer:} B \\
\textbf{Explanation:} This illustrates that dominance depends on the level at which the phenotype is examined.

\item \textbf{Polydactyly (having extra fingers or toes) is caused by a \underline{\hspace{1cm}} allele that is \underline{\hspace{1cm}} in the population.}
\begin{enumerate}[label=\Alph*.]
\item Recessive; rare
\item Dominant; rare
\item Recessive; common
\item Dominant; common
\item Codominant; rare
\end{enumerate}
\textbf{Answer:} B \\
\textbf{Explanation:} Being dominant does not mean an allele is common; most people are homozygous recessive (5 fingers).

\item \textbf{The "norm of reaction" refers to:}
\begin{enumerate}[label=\Alph*.]
\item The speed of a genetic mutation.
\item The phenotypic range of a genotype influenced by the environment.
\item The probability of a cross.
\item The number of alleles for a gene.
\item The distance between two genes.
\end{enumerate}
\textbf{Answer:} B \\
\textbf{Explanation:} For example, hydrangea flower color depends on soil acidity despite the plant's genotype.

\item \textbf{A quantitative character usually indicates:}
\begin{enumerate}[label=\Alph*.]
\item A single gene with two alleles.
\item Polygenic inheritance.
\item Complete dominance.
\item Epistasis only.
\item Environmental influence only.
\end{enumerate}
\textbf{Answer:} B \\
\textbf{Explanation:} Characters that vary continuously (like height) are usually controlled by many genes.

\item \textbf{The dihybrid cross $AaBb \times AaBb$ results in how many different \underline{genotypes}?}
\begin{enumerate}[label=\Alph*.]
\item 4
\item 9
\item 16
\item 8
\item 12
\end{enumerate}
\textbf{Answer:} B \\
\textbf{Explanation:} There are 4 phenotypes but 9 distinct genotypes ($3 \times 3$).

\item \textbf{Mendel’s results were not immediately accepted because they contradicted the popular \underline{\hspace{1cm}} theory of inheritance.}
\begin{enumerate}[label=\Alph*.]
\item Particulate
\item Chromosomal
\item Blending
\item Natural Selection
\item Spontaneous Generation
\end{enumerate}
\textbf{Answer:} C \\
\textbf{Explanation:} People thought traits mixed like paint; Mendel showed they remain discrete.

\item \textbf{What phenotypic ratio is expected in a cross between a dihybrid ($AaBb$) and a double recessive ($aabb$)?}
\begin{enumerate}[label=\Alph*.]
\item 9:3:3:1
\item 1:1:1:1
\item 3:1
\item 1:2:1
\item 9:7
\end{enumerate}
\textbf{Answer:} B \\
\textbf{Explanation:} This is a dihybrid testcross. Each allele combination in the dihybrid gametes appears with equal frequency.

\item \textbf{In a pedigree, a square represents a \underline{\hspace{1cm}} and a circle represents a \underline{\hspace{1cm}}.}
\begin{enumerate}[label=\Alph*.]
\item Female; Male
\item Male; Female
\item Affected; Unaffected
\item Dominant; Recessive
\item Parent; Child
\end{enumerate}
\textbf{Answer:} B \\
\textbf{Explanation:} standard pedigree notation.

\item \textbf{If a man with type AB blood marries a woman with type O blood, what are the possible blood types of their children?}
\begin{enumerate}[label=\Alph*.]
\item A, B, AB, O
\item AB and O only
\item A and B only
\item A, B, O only
\item AB only
\end{enumerate}
\textbf{Answer:} C \\
\textbf{Explanation:} $I^A I^B \times ii$ gives $I^A i$ (Type A) and $I^B i$ (Type B).

\item \textbf{Achondroplasia, a form of dwarfism, is inherited as an autosomal \underline{\hspace{1cm}} trait.}
\begin{enumerate}[label=\Alph*.]
\item Recessive
\item Dominant
\item Codominant
\item Sex-linked
\item Lethal (in heterozygotes)
\end{enumerate}
\textbf{Answer:} B \\
\textbf{Explanation:} Most people are $dd$ (normal height); the dwarf phenotype is $Dd$. $DD$ is lethal.

\item \textbf{In snapdragons, red ($R$) is incompletely dominant over white ($r$). What is the probability of a pink offspring from a $Rr \times rr$ cross?}
\begin{enumerate}[label=\Alph*.]
\item 0\%
\item 25\%
\item 50\%
\item 75\%
\item 100\%
\end{enumerate}
\textbf{Answer:} C \\
\textbf{Explanation:} $Rr$ (pink) and $rr$ (white) are produced in a 1:1 ratio.

\item \textbf{Consanguineous matings (between close relatives) increase the probability of:}
\begin{enumerate}[label=\Alph*.]
\item Dominant mutations.
\item New alleles forming.
\item Homozygosity for rare recessive alleles.
\item Faster evolution.
\item Hybrid vigor.
\end{enumerate}
\textbf{Answer:} C \\
\textbf{Explanation:} Relatives are more likely to carry the same rare recessive genes inherited from a common ancestor.

\item \textbf{What are the possible genotypes for someone with Type A blood?}
\begin{enumerate}[label=\Alph*.]
\item $I^A I^A$ only
\item $I^A i$ only
\item $I^A I^A$ and $I^A i$
\item $I^A I^B$
\item $ii$
\end{enumerate}
\textbf{Answer:} C \\
\textbf{Explanation:} $I^A$ is completely dominant over $i$.

\item \textbf{If a trait is autosomal dominant, an affected child must have:}
\begin{enumerate}[label=\Alph*.]
\item Two affected parents.
\item No affected parents.
\item At least one affected parent.
\item An affected sibling.
\item A carrier parent.
\end{enumerate}
\textbf{Answer:} C \\
\textbf{Explanation:} Dominant traits do not skip generations (unless it's a new mutation).

\item \textbf{If a trait is autosomal recessive, two affected parents will have:}
\begin{enumerate}[label=\Alph*.]
\item Only affected children.
\item 50\% affected children.
\item No affected children.
\item 75\% affected children.
\item Carriers only.
\end{enumerate}
\textbf{Answer:} A \\
\textbf{Explanation:} $aa \times aa$ can only produce $aa$ offspring.

\item \textbf{Albinism is a recessive trait. If two normal parents have an albino child, the parents must be:}
\begin{enumerate}[label=\Alph*.]
\item Homozygous dominant
\item Heterozygous
\item Homozygous recessive
\item Mutants
\item From different species
\end{enumerate}
\textbf{Answer:} B \\
\textbf{Explanation:} Both must be carriers ($Aa$) to contribute an '$a$' allele to the child.

\item \textbf{Cystic fibrosis causes the buildup of mucus in organs because of a defective \underline{\hspace{1cm}} transport protein.}
\begin{enumerate}[label=\Alph*.]
\item Oxygen
\item Chloride
\item Sodium
\item Potassium
\item Glucose
\end{enumerate}
\textbf{Answer:} B \\
\textbf{Explanation:} The CFTR protein failure leads to high extracellular chloride levels and thick mucus.

\item \textbf{A cross between $AABbCc \times AaBbCc$ will produce how many different phenotypes?}
\begin{enumerate}[label=\Alph*.]
\item 2
\item 4
\item 8
\item 1
\item 6
\end{enumerate}
\textbf{Answer:} B \\
\textbf{Explanation:} $A \times Aa$ = 1 phenotype ($A$). $Bb \times Bb$ = 2 phenotypes. $Cc \times Cc$ = 2 phenotypes. $1 \times 2 \times 2 = 4$.

\item \textbf{Which of these is a non-invasive fetal screening tool?}
\begin{enumerate}[label=\Alph*.]
\item Amniocentesis
\item CVS
\item Ultrasound
\item Fetoscopy
\item Bone marrow biopsy
\end{enumerate}
\textbf{Answer:} C \\
\textbf{Explanation:} Ultrasound uses sound waves and does not enter the amniotic sac.

\item \textbf{The phenotype of the heterozygote is identical to the homozygous dominant in:}
\begin{enumerate}[label=\Alph*.]
\item Incomplete dominance
\item Codominance
\item Complete dominance
\item Epistasis
\item Polygenic inheritance
\end{enumerate}
\textbf{Answer:} C \\
\textbf{Explanation:} One allele completely masks the other.

\item \textbf{The dihybrid ratio 9:3:3:1 becomes 9:3:4 in some cases of \underline{\hspace{1cm}}.}
\begin{enumerate}[label=\Alph*.]
\item Pleiotropy
\item Epistasis
\item Incomplete dominance
\item Codominance
\item Linkage
\end{enumerate}
\textbf{Answer:} B \\
\textbf{Explanation:} The 4 results from the combination of the double recessive and one of the single recessives being masked by an upstream gene.

\item \textbf{True-breeding plants are always:}
\begin{enumerate}[label=\Alph*.]
\item Heterozygous
\item Homozygous
\item Recessive
\item Dominant
\item Large
\end{enumerate}
\textbf{Answer:} B \\
\textbf{Explanation:} Only homozygotes produce identical offspring upon selfing.

\item \textbf{Which allele is known as the "null" allele in the ABO system?}
\begin{enumerate}[label=\Alph*.]
\item $I^A$
\item $I^B$
\item $i$
\item All are null
\item None are null
\end{enumerate}
\textbf{Answer:} C \\
\textbf{Explanation:} The '$i$' allele does not produce a functional carbohydrate on the cell surface.

\item \textbf{Mendel’s "Particulate" hypothesis suggests that:}
\begin{enumerate}[label=\Alph*.]
\item Traits mix like liquids.
\item Parents pass on discrete heritable units that retain their identity.
\item Genes are made of particles of dust.
\item Only the father passes on traits.
\item Embryos are pre-formed.
\end{enumerate}
\textbf{Answer:} B \\
\textbf{Explanation:} This was his revolutionary idea that replaced the blending model.

\item \textbf{Probability values range from \underline{\hspace{1cm}} to \underline{\hspace{1cm}}.}
\begin{enumerate}[label=\Alph*.]
\item -1; 1
\item 0; 100
\item 0; 1
\item 1; 10
\item 0; 10
\end{enumerate}
\textbf{Answer:} C \\
\textbf{Explanation:} 0 is certainty of non-occurrence; 1 is certainty of occurrence.

\item \textbf{If $P(A) = 0.4$ and $P(B) = 0.5$, and they are independent, what is $P(A \text{ and } B)$?}
\begin{enumerate}[label=\Alph*.]
\item 0.9
\item 0.1
\item 0.2
\item 0.02
\item 0.5
\end{enumerate}
\textbf{Answer:} C \\
\textbf{Explanation:} $0.4 \times 0.5 = 0.2$.

\item \textbf{How many alleles for a single gene can exist in a \underline{population}?}
\begin{enumerate}[label=\Alph*.]
\item Only 2
\item Only 1
\item Many more than 2
\item Exactly 4
\item 46
\end{enumerate}
\textbf{Answer:} C \\
\textbf{Explanation:} While an individual has 2, a population can have dozens of different alleles for the same locus.

\item \textbf{The substitution in Sickle-cell disease occurs in the \underline{\hspace{1cm}} chain of hemoglobin.}
\begin{enumerate}[label=\Alph*.]
\item Alpha
\item Beta
\item Gamma
\item Delta
\item Epsilon
\end{enumerate}
\textbf{Answer:} B \\
\textbf{Explanation:} It is a specific mutation in the $\beta$-globin gene.

\item \textbf{Genetic testing for adults is often used to identify \underline{\hspace{1cm}}.}
\begin{enumerate}[label=\Alph*.]
\item Their age
\item Carriers of recessive alleles
\item If they are homozygous recessive for all genes
\item Their blood type only
\item Their height
\end{enumerate}
\textbf{Answer:} B \\
\textbf{Explanation:} Prospective parents often test for CF or Tay-Sachs carrier status.

\item \textbf{Mendel’s F2 generation for a monohybrid cross had a genotypic ratio of:}
\begin{enumerate}[label=\Alph*.]
\item 3:1
\item 1:2:1
\item 9:3:3:1
\item 1:1
\item 1:1:1:1
\end{enumerate}
\textbf{Answer:} B \\
\textbf{Explanation:} 1 $PP$ : 2 $Pp$ : 1 $pp$.

\item \textbf{Which character studied by Mendel was dominant in its "inflated" form?}
\begin{enumerate}[label=\Alph*.]
\item Seed shape
\item Pod shape
\item Flower position
\item Stem length
\item Pod color
\end{enumerate}
\textbf{Answer:} B \\
\textbf{Explanation:} Inflated pods were dominant over constricted pods.

\item \textbf{The term "Hybridization" refers to:}
\begin{enumerate}[label=\Alph*.]
\item Self-pollination.
\item Crossing two true-breeding varieties.
\item Cloning an organism.
\item Killing a plant.
\item DNA replication.
\end{enumerate}
\textbf{Answer:} B \\
\textbf{Explanation:} Mendel cross-pollinated different varieties to create hybrids.

\item \textbf{If a homozygous dominant is crossed with a heterozygote, what percentage of offspring will show the recessive phenotype?}
\begin{enumerate}[label=\Alph*.]
\item 0\%
\item 25\%
\item 50\%
\item 75\%
\item 100\%
\end{enumerate}
\textbf{Answer:} A \\
\textbf{Explanation:} $AA \times Aa$ produces only $AA$ and $Aa$ (all dominant phenotype).

\item \textbf{A gene that affects both the heart and the lungs is showing:}
\begin{enumerate}[label=\Alph*.]
\item Epistasis
\item Pleiotropy
\item Polygenic inheritance
\item Incomplete dominance
\item Complete dominance
\end{enumerate}
\textbf{Answer:} B \\
\textbf{Explanation:} One gene, multiple phenotypic effects.

\item \textbf{Which of these is NOT a Mendelian Law?}
\begin{enumerate}[label=\Alph*.]
\item Segregation
\item Independent Assortment
\item Natural Selection
\item All are Mendelian laws
\item None are Mendelian laws
\end{enumerate}
\textbf{Answer:} C \\
\textbf{Explanation:} Natural Selection is Darwin’s principle, not Mendel’s.

\item \textbf{The offspring of the F1 generation are called the \underline{\hspace{1cm}} generation.}
\begin{enumerate}[label=\Alph*.]
\item P
\item F2
\item F3
\item Hybrid
\item Parental
\end{enumerate}
\textbf{Answer:} B \\
\textbf{Explanation:} F2 stands for "second filial."

\item \textbf{Dwarfism from achondroplasia is \underline{\hspace{1cm}} over normal height.}
\begin{enumerate}[label=\Alph*.]
\item Recessive
\item Dominant
\item Codominant
\item Incomplete
\item Epistatic
\end{enumerate}
\textbf{Answer:} B \\
\textbf{Explanation:} The $D$ allele for dwarfism is dominant; $dd$ individuals are normal height.

\item \textbf{In a cross $AaBb \times aabb$, the probability of an $AaBb$ offspring is:}
\begin{enumerate}[label=\Alph*.]
\item 1/16
\item 1/8
\item 1/4
\item 1/2
\item 1
\end{enumerate}
\textbf{Answer:} C \\
\textbf{Explanation:} $Aa \times aa = 1/2 Aa$; $Bb \times bb = 1/2 Bb$. $1/2 \times 1/2 = 1/4$.

\item \textbf{Mendel’s work was "rediscovered" around the year:}
\begin{enumerate}[label=\Alph*.]
\item 1866
\item 1900
\item 1800
\item 1953
\item 1750
\end{enumerate}
\textbf{Answer:} B \\
\textbf{Explanation:} Mendel published in 1866, but his work was ignored until 1900.

\end{enumerate}
\end{document}